\chapter{Some Results from Linear Representation Theory}\label{chapter:appendix}

\section{Induction and Inflation Between Cyclic Groups}

Let $n \geq 1$. Let $\sR_{\bZ}(C_n)$ denote the representation ring of $\bC$-linear representations of $C_n$ over $\bZ$, and let $\sR_{\bC}(C_n)$ denote the base change to $\bC$. For any $m \mid n$ there is an inflation ring homomorphism
\[
\Inf_m^n \colon \sR_{\bZ}(C_m) \rightarrow \sR_{\bZ}(C_n),
\]
and a $\bZ$-linear induction map
\[
\Ind_m^n \colon \sR_{\bZ}(C_m) \rightarrow \sR_{\bZ}(C_n).
\]

Let $Q_n \subset \sR_{\bZ}(C_n)$ denote the image of $\Inf_m^n$ over all $m \mid n$ with $m < n$, which consists of the $\bZ$-span of the non-primitive characters of $C_n$. Let $S_n \subset \sR_{\bZ}(C_n)$ denote the image of $\Inf_m^n$ and $\Ind_m^n$ over all $m \mid n$ with $m < n$.

The aim of this section is to prove the following.

\begin{thm}\label{thm:wheniscyclicnotneeded}
    Suppose that $n \geq 1$. Then $S_n = \sR_{\bZ}(C_n)$ if and only if $n \equiv 2 \! \! \mod 4$.
\end{thm}

The equality when $n \equiv 2 \! \! \mod 4$ is relatively straightforward, but for the converse we will need the following. To state it, for $n \geq 1$ let
\[
\Delta(n) \coloneqq \rank_{\bZ} \sR_{\bZ}(C_n) - \rank_{\bZ} S_n \geq 0.
\]
Let $\mu \colon \bN_{\geq 1} \rightarrow \{-1,0,1\}$ denote the M\"{o}bius function, and $\varphi \colon \bN_{\geq 1} \rightarrow \bN_{\geq 1}$ Euler's totient function.

\begin{prop}\label{prop:descofrankdifference}
    For any $n \geq 1$,
    \[
        \Delta(n) = \sum_{d \mid n} \mu(d) \, \varphi \! \left( \frac{n}{d} \right).
    \]
    In particular, $\Delta \colon \bN_{\geq 1} \rightarrow \bN_{\geq 0}$ is multiplicative, and $\Delta(n) = 0$ if and only if $n \equiv 2 \! \! \mod 4$.
\end{prop}

To prove Proposition \ref{prop:descofrankdifference}, we first show that $\Delta$ is multiplicative.

\begin{lemma}\label{lem:deltaismultiplicitive}
    Suppose that $m,n \geq 1$ with $(m,n) = 1$. Then $\Delta(mn) = \Delta(m) \Delta(n)$.
\end{lemma}

\begin{proof}
    Because $(m,n) = 1$, we have a natural isomorphism
    \[
    \Phi_{m,n} \colon \sR_{\bC}(C_m) \otimes_{\bC} \sR_{\bC}(C_n) \xrightarrow{\sim} \sR_{\bC}(C_m \times C_n).
    \]
    We have that $\Delta(mn)$ is the dimension of the quotient of $\sR_{\bC}(C_m \times C_n)$ by $S_{m,n}$, where $S_{m,n}$ is the sum of the images of the induction maps from proper subgroups, and the inflation maps from proper quotients. By the transitivity of induction and inflation, we may restrict to maximal subgroups, and quotients by minimal subgroups. These are all of the form $C_{m/p} \times C_n$ or $C_m \times C_{n/p}$, where $p$ is a prime divisor of $m$ or $n$. For a prime divisor $p$ of $mn$, if $p \mid m$ there are commutative diagrams
  % https://q.uiver.app/#q=WzAsOCxbMSwwLCJcXHNSX3tcXGJDfShDX20gXFx0aW1lcyBDX24pIl0sWzAsMCwiXFxzUl97XFxiQ30oQ19tKSBcXG90aW1lc197XFxiQ30gXFxzUl97XFxiQ30oQ19uKSJdLFswLDEsIlxcc1Jfe1xcYkN9KENfe20vcH0pIFxcb3RpbWVzX3tcXGJDfSBcXHNSX3tcXGJDfShDX24pIl0sWzEsMSwiXFxzUl97XFxiQ30oQ197bS9wfSBcXHRpbWVzIENfbikiXSxbMiwwLCJcXHNSX3tcXGJDfShDX20pIFxcb3RpbWVzX3tcXGJDfSBcXHNSX3tcXGJDfShDX24pIl0sWzMsMCwiXFxzUl97XFxiQ30oQ19tIFxcdGltZXMgQ19uKSJdLFsyLDEsIlxcc1Jfe1xcYkN9KENfe20vcH0pIFxcb3RpbWVzX3tcXGJDfSBcXHNSX3tcXGJDfShDX24pIl0sWzMsMSwiXFxzUl97XFxiQ30oQ197bS9wfSBcXHRpbWVzIENfbikiXSxbMSwwLCJcXFBoaV97bSxufSIsMl0sWzIsMSwiXFx0aWxkZXtcXEluZH1fe20vcH1ebSBcXG90aW1lcyAxIl0sWzIsMywiXFxQaGlfe20vcCwgbn0iLDJdLFszLDAsIlxcdGlsZGV7XFxJbmR9X3ttL3AsIG59XnttLG59IiwyXSxbMywyLCJcXHNpbSIsMix7InN0eWxlIjp7ImhlYWQiOnsibmFtZSI6Im5vbmUifX19XSxbMCwxLCJcXHNpbSIsMix7InN0eWxlIjp7ImhlYWQiOnsibmFtZSI6Im5vbmUifX19XSxbNCw1LCJcXFBoaV97bSxufSIsMl0sWzYsNywiXFxQaGlfe20vcCxufSIsMl0sWzUsNCwiXFxzaW0iLDIseyJzdHlsZSI6eyJoZWFkIjp7Im5hbWUiOiJub25lIn19fV0sWzcsNiwiXFxzaW0iLDIseyJzdHlsZSI6eyJoZWFkIjp7Im5hbWUiOiJub25lIn19fV0sWzcsNSwiXFx0aWxkZXtcXEluZn1fe20vcCwgbn1ee20sbn0iLDJdLFs2LDQsIlxcdGlsZGV7XFxJbmZ9X3ttL3B9Xm0gXFxvdGltZXMgMSJdXQ==
\[\begin{tikzcd}
	{\sR_{\bC}(C_m) \otimes_{\bC} \sR_{\bC}(C_n)} & {\sR_{\bC}(C_m \times C_n)} & {\sR_{\bC}(C_m) \otimes_{\bC} \sR_{\bC}(C_n)} & {\sR_{\bC}(C_m \times C_n)} \\
	{\sR_{\bC}(C_{m/p}) \otimes_{\bC} \sR_{\bC}(C_n)} & {\sR_{\bC}(C_{m/p} \times C_n)} & {\sR_{\bC}(C_{m/p}) \otimes_{\bC} \sR_{\bC}(C_n)} & {\sR_{\bC}(C_{m/p} \times C_n)}
	\arrow["{\Phi_{m,n}}"', from=1-1, to=1-2]
	\arrow["\sim"', no head, from=1-2, to=1-1]
	\arrow["{\Phi_{m,n}}"', from=1-3, to=1-4]
	\arrow["\sim"', no head, from=1-4, to=1-3]
	\arrow["{\tilde{\Ind}_{m/p}^m \otimes 1}", from=2-1, to=1-1]
	\arrow["{\Phi_{m/p, n}}"', from=2-1, to=2-2]
	\arrow["{\tilde{\Ind}_{m/p, n}^{m,n}}"', from=2-2, to=1-2]
	\arrow["\sim"', no head, from=2-2, to=2-1]
	\arrow["{\tilde{\Inf}_{m/p}^m \otimes 1}", from=2-3, to=1-3]
	\arrow["{\Phi_{m/p,n}}"', from=2-3, to=2-4]
	\arrow["{\tilde{\Inf}_{m/p, n}^{m,n}}"', from=2-4, to=1-4]
	\arrow["\sim"', no head, from=2-4, to=2-3]
\end{tikzcd}\]
and similarly if instead $p \mid n$. In particular, we have that
\[
\Phi_{m,n}(S_m \otimes_{\bC} \sR_{\bC}(C_n) + \sR_{\bC}(C_m) \otimes_{\bC} S_n) = S_{m,n},
\]
and therefore that $\Phi_{m,n}$ induces an isomorphism
\[
    \tilde{\Phi}_{m,n} \colon (\sR_{\bC}(C_m) / S_m) \otimes_{\bC} (\sR_{\bC}(C_n) / S_n) \xrightarrow{\sim} \sR_{\bC}(C_m \times C_n)/ S_{n,m}.
\]
Taking dimensions we see that $\Delta(mn) = \Delta(m) \Delta(n)$ as required.
\end{proof}

\begin{proof}[Proof of Proposition \ref{prop:descofrankdifference}]
    We want to show that $\Delta = \mu * \varphi$, the Dirichlet convolution of $\mu$ and $\varphi$. Because both $\mu$ and $\varphi$ are multiplicative, $\mu * \varphi$ is too, and therefore it suffices to show that $\Delta = \mu * \varphi$ on prime powers. For $p$ a prime and $r \geq 1$, 
    \[
        (\mu * \varphi)(p^r) = \sum_{k = 0}^r \mu(p^k) \varphi(p^{r-k}) = \varphi(p^r) - \varphi(p^{r-1}).
    \]
    On the other hand, $\Delta(p^r)$ is the dimension of the cokernel of
    \[
        \Ind_{p^{r-1}}^{p^r} \oplus \Inf_{p^{r-1}}^{p^r} \colon \sR_{\bC}(C_{p^{r-1}}) \oplus \sR_{\bC}(C_{p^{r-1}}) \rightarrow \sR_{\bC}(C_{p^r}).
    \]
    The image of the inflation map has dimension $p^r - \varphi(p^r)$, consisting of the imprimitive characters, and the image of the induction map has dimension $p^{r-1}$, being injective. The intersection of the images is the image of the non-primitive characters of $C_{p^{r-1}}$, which has dimension $p^{r-1} - \varphi(p^{r-1})$. We may therefore compute that
    \begin{align*}
        \Delta(p^r) = p^r - (p^r - \varphi(p^r)) - p^{r-1} + (p^{r-1} - \varphi(p^{r-1})) = \varphi(p^r) - \varphi(p^{r-1}),
    \end{align*}
    and so $\Delta = \mu * \varphi$, as required. Using this, we have 
    \[
        \Delta(p^r) = \varphi(p^r) - \varphi(p^{r-1}) = \twopartdef{(p-1) - 1}{r = 1}{p^{r-2}(p-1)^2}{r \geq 2},
    \]
    which equals $0$ if and only if $p = 2$ and $r = 1$. In particular, if $n = 2^e p_1^{e_1} \cdots p_k^{e_k}$, then $\Delta(n) = 0$ if and only if $e = 1$, which is the case exactly when $n \equiv 2 \! \! \mod 4$.
\end{proof}

\begin{remark}
    There is a more conceptual proof of Proposition \ref{prop:descofrankdifference} which explains why the M\"{o}bius function appears, at the cost of being much longer than what we have given above.
\end{remark}

\begin{proof}[Proof of Theorem \ref{thm:wheniscyclicnotneeded}]
    When $S_n = \sR_{\bZ}(C_n)$, $\Delta(n) = 0$, and so $n \equiv 2 \! \! \mod 4$ by Proposition \ref{prop:descofrankdifference}. Conversely, suppose that $n \equiv 2 \! \! \mod 4$, so $n = 2m$ for some odd $m \geq 1$. Let $x$ be a generator of $C_n$, let $\chi$ be an irreducible complex character of $C_n$, and set $\zeta \coloneqq \chi(x)$, and $n$th root of $1$. If $\zeta$ is not primitive, then $\chi \in Q_n \subset S_n$. If instead $\zeta$ is primitive, then we can write
    \[
        \chi = \Ind_{m}^n (\psi_{2}) - \Inf_{m}^{n}(\psi_{1+m}),
    \]
    where for $\langle y \rangle = C_m \leq C_n$, $y = x^2$, $\psi_{2}(y) = \zeta^2$ and $\psi_{1+m} = \zeta^{1+m}$, both of which are well-defined characters of $C_m$ because $2$ and $1+m$ are even. Then as $\chi \in S_n$ for any irreducible $\chi$, $S_n = \sR_{\bZ}(C_n)$.
\end{proof}